\documentclass[master.tex]{subfiles}

\begin{document}

% ======================== TITLE ================================
\ifSubfilesClassLoaded{%
  % Standalone title block
  \begin{center}
    {\LARGE\bfseries\color{isublue}
     Chapter 3: The Do-Calculus and Identification Criteria}\\[6pt]
    {\large Theory and Methods in Causal Inference}\\[4pt]
    {\normalsize Department of Statistics, Iowa State University}
  \end{center}
  \bigskip
  \hrule height 1.2pt \relax
  \smallskip
}{%
  % Master book: chapter heading
  \chapter{The Do-Calculus and Identification Criteria}
}
% ======================== LEARNING OBJECTIVES ==================
\begin{tcolorbox}[defstyle, title={Learning Objectives}]
By the end of this chapter, students should be able to:
\begin{enumerate}
  \item Define the two intervention graph operations $\Gcal_{\overline{X}}$ and
        $\Gcal_{\underline{X}}$ and explain their causal interpretation.
  \item Apply the back-door criterion to determine whether a set $\mathbf{S}$
        identifies a causal effect, and write the back-door adjustment formula.
  \item Apply the front-door criterion in graphs where the confounder is
        unobserved but a suitable mediator exists.
  \item State all three rules of do-calculus, identify the graph condition
        that licenses each rule, and use them to prove the back-door and
        front-door formulas.
  \item State the completeness theorem of \citet{shpitser2006identification} and explain
        its practical implication: if the do-calculus cannot identify a quantity,
        no method can.
\end{enumerate}
\end{tcolorbox}

% ---------------------------------------------------------------
% Styles and macros specific to this chapter (not yet in master).
% Placed here (after the LO box, before §1) so the definition block
% does not inject spacing between \chapter{} and the box.
% ---------------------------------------------------------------
% snode: warning-red observed node (used in §2 diagram)
\tikzset{
  snode/.style = {circle, draw=warnred, fill=warnbg, thick,
                  minimum size=7mm, font=\small\bfseries}
}
% Ancestor / Descendant set shorthands (\An, \De, \Nd defined in master.tex)
% Expectation shorthand (\E defined in master.tex)
% Remark style: thin left rule, lighter than defstyle
% (remarkstyle and remark environment are defined in causal_inference_master.tex)
% ==============================================================
\section{From d-Separation to Intervention: Intervention Graphs}
\label{w3:sec:mutilated}
% ==============================================================

In Chapter~2 we learned to read conditional independence from a DAG using
d-separation.  This chapter takes the next step: translating the graphical
language into \emph{identification formulas} --- expressions that write the
interventional distribution $P(y \mid \doop(T{=}t))$ entirely in terms of
quantities observable from data.

\subsection{What Does Intervention Mean?}
\label{w3:sec:intervention-meaning}

Before introducing any graph machinery, it is worth being precise about what
the $\doop(\cdot)$ operator means and why it is not the same as conditioning.

\paragraph{Conditioning versus intervening.}
The conditional distribution $P(Y \mid T{=}t)$ describes the subpopulation of
units for whom $T$ was \emph{observed} to equal $t$.  Because treatment
assignment may be influenced by confounders, this subpopulation is not
representative of the full population.  The interventional distribution
$P(Y \mid \doop(T{=}t))$, by contrast, describes the population that would
result if $T$ were \emph{externally set} to $t$ for everyone --- removing
the natural process that determines $T$ and replacing it with the forced
value.  The two distributions coincide only when $T$ has no unmeasured
common causes with $Y$.

A simple illustration: suppose temperature $Z$ causes both ice-cream sales
$T$ and crime $Y$.
\begin{itemize}
  \item $P(Y \mid T{=}t)$ is the crime rate on days when ice-cream sales
        happen to equal $t$ --- high sales happen on hot days, which also
        drive higher crime, so the correlation is confounded by $Z$.
  \item $P(Y \mid \doop(T{=}t))$ is the crime rate we would observe if we
        \emph{fixed} sales at level $t$ by decree --- say, by regulating
        prices.  Forcing sales does not change the temperature, so crime
        remains driven only by $Z$, not by the sales level we mandated.
        If $T$ has no direct causal path to $Y$, then
        $P(Y \mid \doop(T{=}t)) = P(Y)$ for all $t$.
\end{itemize}

\paragraph{The structural basis of the do-operator.}
In the SEM framework of Chapter~1, every variable is generated by a
structural equation.  For the treatment node:
\[
  T \;=\; f_T\!\bigl(\Pa(T),\, U_T\bigr),
\]
where $\Pa(T)$ are the causal parents of $T$ and $U_T$ is exogenous noise.
An \emph{intervention} $\doop(T{=}t)$ \textbf{replaces this entire equation}
with the constant $T = t$.  The replacement has two consequences:
\begin{enumerate}
  \item $T$ is no longer influenced by its parents --- its dependence on
        $\Pa(T)$ and $U_T$ is severed.
  \item All other structural equations remain unchanged --- the rest of the
        causal mechanism is unaffected.
\end{enumerate}
This is the sense in which an intervention is a \emph{surgical} operation on
the model: it cuts out one equation and replaces it with a constant, leaving
the remainder of the system intact.

\paragraph{Graph surgery as the graphical realization.}
Because $\Pa(T)$ no longer affects $T$ after the intervention, every arrow
pointing \emph{into} $T$ in the DAG becomes meaningless and should be
removed.  The resulting graph --- the \emph{intervention graph}
$\Gcal_{\overline{T}}$ --- represents the post-intervention world.
D-separation in $\Gcal_{\overline{T}}$ therefore encodes conditional
independence in $P(\,\cdot \mid \doop(T{=}t))$, not in the original
observational distribution $P$.  This is the key link that allows
graphical reasoning to answer causal questions.

This chapter follows a \emph{concrete-to-general} order.
Sections~\ref{w3:sec:backdoor} and~\ref{w3:sec:frontdoor}
present the two most important identification strategies --- the back-door and
front-door criteria --- as self-contained results, building intuition through
worked examples.  Section~\ref{w3:sec:rules} then reveals the three rules of
do-calculus as the general algebraic engine that underpins both, and supplies
their proofs.

\subsection{The Two Graph Operations}
\label{w3:sec:operations}
Both operations modify the DAG by selectively removing edges.

\begin{defn}{Intervention Graphs $\Gcal_{\overline{X}}$ and $\Gcal_{\underline{X}}$ (Pearl's ``Mutilated Graphs'')}
Let $\Gcal = (\mathcal{V}, \mathcal{E})$ be a DAG and $X \subseteq
\mathcal{V}$ a set of nodes.
\begin{itemize}
  \item $\Gcal_{\overline{X}}$ (\emph{intervention graph}, or
        \emph{intervention graph} in Pearl's terminology) is obtained by
        \emph{deleting all edges into} $X$: remove every arrow
        $V \to X_i$ for all $X_i \in X$.
        This represents the intervention $\doop(X{=}x)$, which
        \emph{replaces the structural equation for $X$ by the constant $x$},
        thereby severing $X$'s dependence on all its former parents.
  \item $\Gcal_{\underline{X}}$ (\emph{auxiliary observation graph}) is
        obtained by \emph{deleting all edges out of} $X$: remove every
        arrow $X_i \to V$ for all $X_i \in X$.
        This is a \emph{technical device} used in Rule~1 of the do-calculus
        to determine when an observation on $X$ can be inserted into or
        removed from a probability expression.  It does \emph{not} represent
        a probabilistic model of conditioning on $X$; it is simply the graph
        in which the relevant d-separation condition for Rule~1 must be
        checked.
\end{itemize}
\end{defn}

\noindent\textbf{Mnemonic.}
\[
  \Gcal_{\overline{X}}: \text{cut \emph{incoming} arrows into } X,
  \qquad
  \Gcal_{\underline{X}}: \text{cut \emph{outgoing} arrows from } X.
\]

Figure~\ref{w3:fig:mutilations} illustrates both intervention graph operations on the IV DAG
$\{Z \to T,\; T \to Y,\; U \to T,\; U \to Y\}$.

\begin{figure}[h]
\centering
\begin{tikzpicture}[node distance=1.8cm]
  %% ---- Original DAG ----
  \node[node]  (Z0) at (0,0)       {$Z$};
  \node[node]  (T0) at (1.9,0)     {$T$};
  \node[node]  (Y0) at (3.8,0)     {$Y$};
  \node[unode] (U0) at (2.85,-1.3) {$U$};
  \draw[edge]  (Z0) -- (T0);
  \draw[edge]  (T0) -- (Y0);
  \draw[dedge] (U0) -- (T0);
  \draw[dedge] (U0) -- (Y0);
  \node[font=\small\itshape\color{darkgrey}] at (1.9,-2.2) {Original $\Gcal$};

  %% ---- G_T-bar: delete edges INTO T ----
  \node[node]  (Z1) at (5.8,0)      {$Z$};
  \node[node]  (T1) at (7.7,0)      {$T$};
  \node[node]  (Y1) at (9.6,0)      {$Y$};
  \node[unode] (U1) at (8.65,-1.3)  {$U$};
  \draw[xedge] (Z1) -- (T1);
  \draw[edge]  (T1) -- (Y1);
  \draw[xedge] (U1) -- (T1);
  \draw[dedge] (U1) -- (Y1);
  \node[font=\small\itshape\color{darkgrey}] at (7.7,-2.2)
    {$\Gcal_{\overline{T}}$: delete arrows \emph{into} $T$};

  %% ---- G_T-underline: delete edges OUT of T ----
  \node[node]  (Z2) at (11.6,0)     {$Z$};
  \node[node]  (T2) at (13.5,0)     {$T$};
  \node[node]  (Y2) at (15.4,0)     {$Y$};
  \node[unode] (U2) at (14.45,-1.3) {$U$};
  \draw[edge]  (Z2) -- (T2);
  \draw[xedge] (T2) -- (Y2);
  \draw[dedge] (U2) -- (T2);
  \draw[dedge] (U2) -- (Y2);
  \node[font=\small\itshape\color{darkgrey}] at (13.5,-2.2)
    {$\Gcal_{\underline{T}}$: delete arrows \emph{out of} $T$};
\end{tikzpicture}
\caption{The two graph operations on the IV DAG.  Faded arrows have been
removed.  $\Gcal_{\overline{T}}$ models external intervention $\doop(T{=}t)$:
the back-door path $T\leftarrow U \to Y$ is severed.
$\Gcal_{\underline{T}}$ deactivates $T$'s influence on its descendants,
used by Rule~1 to add or delete $T$ from a conditioning set.}
\label{w3:fig:mutilations}
\end{figure}

\subsection{Why the Intervention Graph Works}
\label{w3:sec:mutilation-principle}
The causal interpretation rests on a key structural insight:

\begin{tcolorbox}[defstyle, title={The Intervention Principle}]
In the original DAG $\Gcal$, the node $T$ is determined by its parents
$\Pa(T)$ via the structural equation $T = f_T(\Pa(T), \varepsilon_T)$.
Intervening on $T$ removes the structural mechanism that normally determines
$T$; therefore the parents of $T$ no longer influence it.
The intervention $\doop(T{=}t)$ \emph{replaces this equation} with the
constant $T = t$, making $T$ independent of all its former parents.
Graphically, this severs every arrow into $T$.  The remaining structure of
the graph is unchanged: all other structural equations still hold.

A d-separation statement in $\Gcal_{\overline{T}}$ is therefore a
conditional independence statement in the \emph{post-intervention}
distribution $P(\,\cdot \mid \doop(T{=}t))$, not in the observational
distribution $P$.  Section~\ref{w3:sec:rules} develops the three do-calculus
rules that make this correspondence precise and supplies the proofs of the
formulas introduced in Sections~\ref{w3:sec:backdoor}
and~\ref{w3:sec:frontdoor}.
\end{tcolorbox}

% ==============================================================
\section{The Back-Door Criterion}
\label{w3:sec:backdoor}
% ==============================================================

The back-door criterion is the most widely used identification strategy.
It gives a simple graphical condition under which conditioning on an
observed set $\mathbf{S}$ is sufficient to identify the causal effect of
$T$ on $Y$.  The formula it produces --- standardization over $\mathbf{S}$
--- is computable directly from the observed data distribution;
Section~\ref{w3:sec:proofs} supplies the do-calculus proof.

\subsection{Motivation: When the Confounder Is Unobserved}
\label{w3:sec:backdoor-motiv}

If all common causes of $T$ and $Y$ were observed, identification would be trivial: adjust
for all of them.  The interesting case --- and the one the back-door criterion is designed
for --- is when some confounders are \emph{unobserved}, yet an observed set $\mathbf{S}$
can still block every spurious path.

Consider the following DAG.  A drug $T$ affects recovery $Y$.
Socioeconomic status $U$ is an unobserved confounder ($U \to T$, $U \to Y$),
but $U$ also affects an observed proxy $W$ ($U \to W$), which in turn affects $Y$
($W \to Y$).

\begin{center}
\begin{tikzpicture}[node distance=1.8cm]
  \node[unode]  (U) at (1.8,  0)   {$U$};
  \node[node]   (W) at (3.6,  0)   {$W$};
  \node[node]   (T) at (0,   -1.4) {$T$};
  \node[node]   (Y) at (3.6, -1.4) {$Y$};
  \draw[edge]   (T) -- (Y);
  \draw[dedge]  (U) -- (T);
  \draw[dedge]  (U) -- (W);
  \draw[edge]   (W) -- (Y);
\end{tikzpicture}
\end{center}

\noindent
The only back-door path is $T \leftarrow U \to W \to Y$.  Although $U$ is unobserved
and cannot be adjusted for directly, $W$ sits on this path and \emph{is} observed.
Conditioning on $W$ blocks the path at the chain node $W$, and $W$ is not a descendant
of $T$, so $\mathbf{S} = \{W\}$ satisfies the back-door criterion.  The causal effect is
therefore identified even though $U$ is never measured.

This is the key insight: the criterion does not require adjusting for the confounders
themselves, only for an observed set that \emph{intercepts} every back-door path.

\subsection{Definition and Theorem}
\label{w3:sec:backdoor-def}
\begin{defn}{Back-Door Criterion}
\label{w3:def:backdoor}
A set of observed variables $\mathbf{S}$ satisfies the \emph{back-door
criterion} for the causal effect of $T$ on $Y$ in $\Gcal$ if:
\begin{enumerate}
  \item No node in $\mathbf{S}$ is a descendant of $T$.
  \item $\mathbf{S}$ blocks every \emph{back-door path} from $T$ to $Y$,
        i.e., every path that begins with an arrow pointing \emph{into}
        $T$ (a path of the form $T \leftarrow \cdots Y$).
\end{enumerate}
\end{defn}

\begin{remark}
The terminology reflects the geometry of the graph relative to $T$.
A \emph{back-door path} enters $T$ from behind --- it begins with an arrow
pointing \emph{into} $T$ --- and represents a confounding route through
which $T$ and $Y$ become spuriously associated without any causal mechanism
from $T$ to $Y$.  By contrast, the directed causal paths
$T \to \cdots \to Y$ leave $T$ through the \emph{front} and carry the
genuine effect.  The front-door criterion of
Section~\ref{w3:sec:frontdoor} is named for exactly this complementary
idea: it identifies the causal effect by intercepting the \emph{front-door}
(causal) paths at a mediator $M$, rather than blocking the back-door
(confounding) paths directly.
Standing at $T$, confounders sneak in through the back door; the causal
effect walks out through the front.
\end{remark}

\begin{remark}
The adjustment set $\mathbf{S}$ here plays the same role as the
conditioning set in the d-separation notation of Chapter~2.
Specifically, condition~2 is
equivalent to requiring $(Y \indep T \mid \mathbf{S})_{\Gcal_{\overline{T}}}$:
$\mathbf{S}$ d-separates $T$ and $Y$ in the intervention graph.  The
back-door criterion adds two constraints that $\mathbf{S}$ in Chapter~2
did not carry: $\mathbf{S}$ must consist of \emph{observed} variables
only, and it must contain no descendant of $T$.  Every valid adjustment
set $\mathbf{S}$ is therefore a valid d-separator, but the converse does
not hold --- a set that d-separates $T$ and $Y$ may still fail the
back-door criterion if it contains a descendant of $T$ or an unobserved
variable.
\end{remark}

\begin{example}{A Single Confounder}
\label{w3:ex:singleconf}
Consider the DAG $T \leftarrow C \to Y$ with $T \to Y$, where $C$ is an
observed confounder.  We wish to identify $P(y \mid \doop(T{=}t))$.

\begin{center}
\begin{tikzpicture}[node distance=1.8cm]
  \node[node] (C) at (1.8, 0)   {$C$};
  \node[node] (T) at (0,  -1.4) {$T$};
  \node[node] (Y) at (3.6,-1.4) {$Y$};
  \draw[edge] (T) -- (Y);
  \draw[edge] (C) -- (T);
  \draw[edge] (C) -- (Y);
\end{tikzpicture}
\end{center}

\noindent
There is one back-door path: $T \leftarrow C \to Y$ (an arrow points into
$T$).  Does $\mathbf{S} = \{C\}$ satisfy the back-door criterion?
\begin{enumerate}
  \item $C$ is not a descendant of $T$. \checkmark
  \item $\{C\}$ blocks $T \leftarrow C \to Y$ (fork at $C$). \checkmark
\end{enumerate}
The criterion is satisfied, and the back-door formula gives:
\[
  P(y \mid \doop(T{=}t)) = \sum_c P(y \mid t, c)\,P(c).
\]
This is the simplest possible application: one confounder, one back-door
path, one adjustment variable.
\end{example}

\begin{warn}{Why Condition~1 Is Essential: A Worked Counterexample}
Condition~1 rules out conditioning on any descendant of $T$.  To see why,
consider a drug $T$ that affects recovery $Y$ via two routes: a direct effect
$T \to Y$ and an indirect effect through blood pressure $M$ ($T \to M \to Y$).
Suppose also that an unobserved confounder $U$ affects both $T$ and $Y$.

\begin{center}
\begin{tikzpicture}[node distance=1.8cm]
  \node[unode] (U) at (1.8,  0)   {$U$};
  \node[node]  (T) at (0,   -1.4) {$T$};
  \node[node]  (M) at (1.8, -1.4) {$M$};
  \node[node]  (Y) at (3.6, -1.4) {$Y$};
  \draw[edge]  (T) -- (M);
  \draw[edge]  (M) -- (Y);
  \draw[edge]  (T) to[bend right=30] (Y);
  \draw[dedge] (U) -- (T);
  \draw[dedge] (U) -- (Y);
\end{tikzpicture}
\end{center}

\noindent
The back-door path is $T \leftarrow U \to Y$.  Now suppose we naively try
$\mathbf{S} = \{M\}$, hoping that conditioning on $M$ will block this path.
It does not: $M$ lies on the \emph{causal} path $T \to M \to Y$, not on the
back-door path $T \leftarrow U \to Y$.  Conditioning on $M$ therefore fails
for two distinct reasons:
\begin{enumerate}
  \item It violates Condition~1: $M$ is a descendant of $T$.
  \item It \emph{blocks causal information}: the stratum-specific association
        $P(y \mid t, m)$ is computed holding $M$ fixed, so it cannot see the
        effect of $T$ that flows through $M$.  The result $\sum_m P(y \mid t,m)P(m)$
        estimates the direct effect $T \to Y$ only, not the total causal effect.
\end{enumerate}
The correct adjustment for the total effect uses an observed set that blocks the
back-door path $T \leftarrow U \to Y$ \emph{without touching} the causal paths.
If $U$ is unobserved and no other observed variable sits on the back-door path,
the total causal effect is \emph{not identified} by back-door adjustment in this graph.
\end{warn}

\begin{thm}{Back-Door Adjustment Formula \citep{pearl1993comment}}
\label{w3:thm:backdoor}
If $\mathbf{S}$ satisfies the back-door criterion for the effect of $T$ on
$Y$ in $\Gcal$, then:
\begin{equation}
  P\!\left(y \mid \doop(T{=}t)\right)
  \;=\; \int P(y \mid T{=}t,\;\mathbf{S}{=}\mathbf{s})\,dP(\mathbf{s}).
  \label{w3:eq:backdoor}
\end{equation}
Equivalently, for any measurable function $h$ with $\E|h(Y)| < \infty$,
\begin{equation}
  \E\!\left[h(Y) \mid \doop(T{=}t)\right]
  \;=\; \E_{\mathbf{S}}\!\Bigl[\E\bigl[h(Y) \mid T{=}t,\;\mathbf{S}\bigr]\Bigr].
  \label{w3:eq:backdoor-exp}
\end{equation}
The interventional distribution is identified by integrating the
conditional distribution $P(y \mid t, \mathbf{s})$ with respect to the
\emph{marginal} distribution of $\mathbf{S}$.  The iterated-expectation
form~\eqref{w3:eq:backdoor-exp} makes this structure transparent: the
outer expectation averages over the population distribution of $\mathbf{S}$,
while the inner expectation computes the stratum-specific association --- which,
within each stratum, is causal because all back-door paths are blocked.
When $\mathbf{S}$ is discrete, the integral reduces to
$\sum_{\mathbf{s}} P(y \mid t, \mathbf{s})\,P(\mathbf{s})$.
\end{thm}

The derivation uses Rules~2 and~3 of the do-calculus; the complete
three-step proof is given in Section~\ref{w3:sec:proofs}.

\subsection{The Adjustment Formula: Standardization}
\label{w3:sec:standardization}
Equation~\eqref{w3:eq:backdoor} is also known as the \emph{standardization
formula} or, in the epidemiological literature, the \emph{g-formula}
\citep{robins1986new}.  The g-formula was originally derived for the single
time-point case $K=1$; its multi-period generalization to dynamic treatment
regimes --- where treatment and covariates are both time-varying --- is the
subject of Chapter~13.  The interpretation of the formula is important:

\begin{itemize}
  \item $P(y \mid t, \mathbf{s})$ is the stratum-specific conditional
        distribution of $Y$ given $T{=}t$ and $\mathbf{S}{=}\mathbf{s}$.
        Within each stratum, $T$ is not confounded by any back-door path
        (all are blocked by $\mathbf{S}$), so this conditional distribution
        \emph{is} causal.
  \item The outer integration $\int \cdots\, dP(\mathbf{s})$ weights the
        strata by the \emph{population marginal} distribution of
        $\mathbf{S}$, rather than the treatment-conditional distribution
        $dP(\mathbf{s} \mid t)$ that a naive analysis would use.  The
        iterated-expectation form~\eqref{w3:eq:backdoor-exp} makes this
        re-weighting explicit: $\E_{\mathbf{S}}[\cdots]$ is expectation
        under $P(\mathbf{S})$, not $P(\mathbf{S} \mid T{=}t)$.
  \item In a randomized experiment, $\mathbf{S} = \varnothing$ satisfies
        the back-door criterion (no confounding), and
        \eqref{w3:eq:backdoor} reduces to $P(y \mid \doop(t)) = P(y \mid t)$:
        the observational and interventional distributions coincide.
\end{itemize}

\subsection{Application to the Education--Earnings DAG}
\label{w3:sec:education-ex}
\begin{example}{Back-Door Identification in the Education Example}
\label{w3:ex:education}
This is the Education--Earnings DAG analyzed in full in
Section~\ref{w2:sec:education} of Chapter~2, where Steps~1--3 of the
identification pipeline were carried out.  We now complete Step~4: applying the back-door criterion to
derive the identification formula.  The DAG has edges $N \to E$, $B \to E$,
$B \to Y$, $E \to Y$, where $N$ = neighborhood, $B$ = family background,
$E$ = education, $Y$ = earnings.  We wish to identify
$P(y \mid \doop(E{=}e))$.

\begin{center}
\begin{tikzpicture}[node distance=1.8cm]
  \node[node] (N) at (0, 0)     {$N$};
  \node[node] (B) at (2.0, 0)   {$B$};
  \node[node] (E) at (1.0,-1.4) {$E$};
  \node[node] (Y) at (3.0,-1.4) {$Y$};
  \draw[edge] (E) -- (Y);
  \draw[edge] (N) -- (E);
  \draw[-{Stealth[length=5pt]}, thick, color=warnred, dashed] (B) -- (E);
  \draw[-{Stealth[length=5pt]}, thick, color=warnred, dashed] (B) -- (Y);
  \node[font=\scriptsize\itshape\color{darkgrey}, above=0.05cm of N] {neighborhood};
  \node[font=\scriptsize\itshape\color{darkgrey}, above=0.05cm of B] {background};
  \node[font=\scriptsize\itshape\color{darkgrey}, below=0.05cm of E] {education};
  \node[font=\scriptsize\itshape\color{darkgrey}, below=0.05cm of Y] {earnings};
  \draw[-{Stealth[length=5pt]}, thick, color=warnred, dashed]
    (0.0,-2.3) -- (0.6,-2.3)
    node[anchor=west, font=\small\color{warnred}]
    {\; back-door path $E \leftarrow B \to Y$};
\end{tikzpicture}
\end{center}

\medskip\noindent
\textbf{Back-door paths from $E$ to $Y$:}
The only back-door path is $E \leftarrow B \to Y$ (an arrow points into
$E$ from $B$, then continues to $Y$).

\medskip\noindent
\textbf{Does $\{B\}$ satisfy the criterion?}
\begin{enumerate}
  \item $B$ is not a descendant of $E$. \checkmark
  \item $\{B\}$ blocks $E \leftarrow B \to Y$ (fork at $B$,
        conditioning closes it). \checkmark
\end{enumerate}
By Theorem~\ref{w3:thm:backdoor}:
\[
  P(y \mid \doop(E{=}e)) = \sum_b P(y \mid e, b)\,P(b).
\]

\medskip\noindent
\textbf{Does $\{N\}$ satisfy the criterion?}
$N$ is not a descendant of $E$, but $\{N\}$ does \emph{not} block
$E \leftarrow B \to Y$ (the fork at $B$ remains open).  So $\{N\}$ fails
condition~2 and cannot be used for adjustment.

\medskip\noindent
\textbf{Does $\{N, B\}$ satisfy the criterion?}
Both conditions hold: neither $N$ nor $B$ is a descendant of $E$, and
$\{N,B\}$ blocks the only back-door path.  So $\{N,B\}$ also yields a
valid (though redundant) adjustment formula.
\end{example}

% ==============================================================
\section{The Front-Door Criterion}
\label{w3:sec:frontdoor}
% ==============================================================

The back-door criterion requires that all confounders be observed.  In
many applications --- including the IV setting --- the confounder $U$ is
unobserved and no back-door adjustment is possible.  The front-door
criterion provides identification in a different configuration: when there
exists a measured \emph{mediator} $M$ on the causal path from $T$ to $Y$
with specific independence properties.

\subsection{The Configuration and Criterion}
\label{w3:sec:frontdoor-def}
\begin{defn}{Front-Door Criterion}
\label{w3:def:frontdoor}
A set of variables $\mathbf{M}$ satisfies the \emph{front-door criterion}
for the effect of $T$ on $Y$ in $\Gcal$ if:
\begin{enumerate}
  \item $\mathbf{M}$ intercepts all directed paths from $T$ to $Y$
        (all causal paths pass through $\mathbf{M}$).
  \item There are no unblocked back-door paths from $T$ to $\mathbf{M}$
        (the effect of $T$ on $\mathbf{M}$ is not confounded).
  \item All back-door paths from $\mathbf{M}$ to $Y$ are blocked by $T$
        (conditioning on $T$ closes all confounding of the
        $\mathbf{M} \to Y$ effect).
\end{enumerate}
\end{defn}

\begin{example}{A Single Mediator with Unobserved Confounding}
\label{w3:ex:frontdoorconds}
Consider the DAG $T \to M \to Y$ with $U \to T$ and $U \to Y$, where
$M$ is observed and $U$ is unobserved.

\begin{center}
\begin{tikzpicture}[node distance=1.8cm]
  \node[node]  (T) at (0,0)      {$T$};
  \node[node]  (M) at (2.0,0)    {$M$};
  \node[node]  (Y) at (4.0,0)    {$Y$};
  \node[unode] (U) at (2.0,-1.4) {$U$};
  \draw[edge]  (T) -- (M);
  \draw[edge]  (M) -- (Y);
  \draw[dedge] (U) -- (T);
  \draw[dedge] (U) -- (Y);
\end{tikzpicture}
\end{center}

\noindent Does $M$ satisfy the front-door criterion?
\begin{enumerate}
  \item \emph{$M$ intercepts all directed paths from $T$ to $Y$.}
        The only directed path is $T \to M \to Y$, which passes through
        $M$. \checkmark
  \item \emph{No unblocked back-door paths from $T$ to $M$.}
        The only path from $T$ to $M$ that is not the direct arrow
        $T \to M$ would have to go through $U$, but there is no edge
        $U \to M$, so no back-door path exists. \checkmark
  \item \emph{All back-door paths from $M$ to $Y$ are blocked by $T$.}
        The only back-door path from $M$ to $Y$ is
        $M \leftarrow T \leftarrow U \to Y$.
        Conditioning on $T$ blocks this path (chain at $T$). \checkmark
\end{enumerate}
All three conditions hold, so the front-door formula applies:
\[
  P(y \mid \doop(T{=}t))
  = \sum_m P(m \mid t)\sum_{t'} P(y \mid t', m)\,P(t').
\]
\end{example}

The prototypical graph where the front-door criterion applies is:
\[
  T \;\to\; M \;\to\; Y, \qquad U \to T, \quad U \to Y,
  \quad \text{no direct } T \to Y \text{ edge.}
\]
The absence of a direct $T \to Y$ edge is not an independent requirement:
it is simply one way of satisfying condition~1, which demands that
\emph{all} directed paths from $T$ to $Y$ pass through $M$.  A direct
$T \to Y$ edge would create a second causal path bypassing $M$, violating
condition~1 and breaking the front-door argument --- as Problem~3(c) below
illustrates.
Here $U$ is an unobserved common cause of $T$ and $Y$; $M$ is a
fully-observed mediator.  Figure~\ref{w3:fig:frontdoor} shows this graph.

\begin{figure}[h]
\centering
\begin{tikzpicture}[node distance=1.8cm]
  \node[node]  (T) at (0,0)      {$T$};
  \node[node]  (M) at (2.0,0)    {$M$};
  \node[node]  (Y) at (4.0,0)    {$Y$};
  \node[unode] (U) at (2.0,-1.4) {$U$};
  \draw[edge]  (T) -- (M);
  \draw[edge]  (M) -- (Y);
  \draw[dedge] (U) -- (T);
  \draw[dedge] (U) -- (Y);
  \node[font=\small\itshape\color{darkgrey}, below=0.15cm of T] {treatment};
  \node[font=\small\itshape\color{darkgrey}, below=0.15cm of M] {mediator};
  \node[font=\small\itshape\color{darkgrey}, below=0.15cm of Y] {outcome};
  \node[font=\small\itshape\color{darkgrey}, below=0.15cm of U] {unobserved};
\end{tikzpicture}
\caption{The prototypical front-door graph.  The confounder $U$ is
unobserved, so back-door adjustment on $U$ is infeasible.  However, the
mediator $M$ is fully observed and satisfies the three front-door
conditions, enabling identification via the front-door formula.}
\label{w3:fig:frontdoor}
\end{figure}

\subsection{The Front-Door Formula}
\label{w3:sec:frontdoor-formula}
\begin{thm}{Front-Door Adjustment Formula \citep{pearl1995causal}}
\label{w3:thm:frontdoor}
If $M$ satisfies the front-door criterion for the effect of $T$ on $Y$, and
$P(t) > 0$ for all $t$, then:
\begin{equation}
  P\!\left(y \mid \doop(T{=}t)\right)
  \;=\; \int
      \underbrace{\Bigl[\int P(y \mid T{=}t',\;M{=}m)\,dP(t')\Bigr]}_{\displaystyle\text{Stage 2}}
      \;\underbrace{dP(m \mid t)}_{\displaystyle\text{Stage 1}}.
  \label{w3:eq:frontdoor}
\end{equation}
Equivalently, writing $\mu(t',m) = \E[Y \mid T{=}t',\;M{=}m]$,
\begin{equation}
  \E\!\left[Y \mid \doop(T{=}t)\right]
  \;=\; \E_{M \mid T=t}\!\Bigl[\E_T\bigl[\mu(T,\,M)\bigr]\Bigr].
  \label{w3:eq:frontdoor-exp}
\end{equation}
\textbf{Stage~1}: $dP(m \mid t)$ is the distribution of the mediator $M$
given $T{=}t$.  Front-door condition~2 guarantees no confounding between
$T$ and $M$, so this conditional distribution is directly identifiable
from the observed data.\\[4pt]
\textbf{Stage~2}: $\int P(y \mid t',\,m)\,dP(t')$ is the back-door-adjusted
causal effect of $M$ on $Y$, with $T$ serving as the adjustment variable.
In the iterated-expectation form~\eqref{w3:eq:frontdoor-exp}, this is
$\E_T[\mu(T,M)]$: the regression function $\mu(T,M)$ averaged over the
\emph{marginal} distribution of $T$, not the conditional distribution
$P(T \mid M{=}m)$.  Front-door condition~3 guarantees that conditioning
on $T$ blocks all back-door paths from $M$ to $Y$ (specifically,
$M \leftarrow T \leftarrow U \to Y$), so this marginalization removes the
confounding introduced by $U$.\\[4pt]
The formula chains the two stages by integrating Stage~2 against the
Stage~1 measure $dP(m \mid t)$.  When $T$ and $M$ are both discrete,
the integrals reduce to the summation form
$\sum_m P(m\mid t)\sum_{t'} P(y\mid t',m)\,P(t')$.
\end{thm}

\noindent\textit{Derivation sketch.}
The formula is the composition of two back-door applications.
\begin{enumerate}
  \item \textbf{Identify $P(m \mid \doop(T{=}t))$.}
        By condition~2 there are no back-door paths from $T$ to $M$,
        so $T$ satisfies the (trivial) back-door criterion for its effect
        on $M$ with empty adjustment set.  Hence
        $P(m \mid \doop(T{=}t)) = P(m \mid T{=}t)$.

  \item \textbf{Identify $P(y \mid \doop(M{=}m))$.}
        By condition~3, every back-door path from $M$ to $Y$ passes
        through $T$.  Conditioning on $T$ blocks those paths, so $T$ is
        a valid back-door adjustment set for the $M \to Y$ effect:
        \[
          P(y \mid \doop(M{=}m))
          \;=\; \int P(y \mid T{=}t',\;M{=}m)\,dP(t').
        \]

  \item \textbf{Chain the two stages.}
        By condition~1, every directed path from $T$ to $Y$ passes
        through $M$, so the effect of $T$ on $Y$ operates entirely
        \emph{through} $M$.  Combining steps~1 and~2:
        \[
          P(y \mid \doop(T{=}t))
          = \int P(y \mid \doop(M{=}m))\,dP(m \mid \doop(T{=}t))
          = \int \Bigl[\int P(y \mid t',m)\,dP(t')\Bigr]\,dP(m \mid t).
        \]
\end{enumerate}
The role of each condition is now transparent: condition~1 licenses
step~3 (the chaining), condition~2 licenses step~1 (replacing
$\doop(T{=}t)$ with conditioning), and condition~3 licenses step~2
(back-door adjustment for $M \to Y$).

A formal proof via three applications of the do-calculus rules is given in
Section~\ref{w3:sec:proofs}.

\begin{figure}[t]
\centering
\begin{tikzpicture}[node distance=1.8cm and 2.0cm]
  \node[node]  (T) at (0,0)     {$T$};
  \node[node]  (M) at (3.2,0)   {$M$};
  \node[node]  (Y) at (6.4,0)   {$Y$};
  \node[unode] (U) at (3.2,2.4) {$U$};

  \draw[edge]  (T) -- (M);
  \draw[edge]  (M) -- (Y);
  \draw[dedge] (U) -- (T);
  \draw[dedge] (U) -- (Y);

  \node[font=\small\itshape\color{darkgrey}, align=center]
    at (1.6,-0.90) {Stage~1\\identify $T \to M$};
  \node[font=\small\itshape\color{darkgrey}, align=center]
    at (4.8,-0.90) {Stage~2\\identify $M \to Y \mid T$};

  \node[font=\small\itshape\color{darkgrey}]
    at (3.2,3.10) {$U$ unobserved};

  \node[font=\scriptsize\itshape\color{darkgrey}, align=center]
    at (3.2,-1.85)
    {(conditioning on $T$ blocks $M \leftarrow T \leftarrow U \to Y$,
      making the $M \to Y$ effect identifiable)};
\end{tikzpicture}
\caption{Why the front-door criterion works.  The total effect of $T$ on $Y$
is not directly identifiable because $U$ confounds $T$ and $Y$.  However,
the effect of $T$ on the mediator $M$ is unconfounded (Stage~1), and the
effect of $M$ on $Y$ becomes identifiable after conditioning on $T$, which
blocks the back-door path through $U$ (Stage~2).  The front-door formula
reconstructs the total causal effect by combining these two identifiable
pieces.}
\label{w3:fig:frontdoor-two-stage}
\end{figure}

\subsection{Interpretation: Two-Stage Causal Accounting}
\label{w3:sec:twostage}
The two-stage structure annotated in equation~\eqref{w3:eq:frontdoor} has a
natural graphical reading in Figure~\ref{w3:fig:frontdoor-two-stage}: Stage~1
traverses the unconfounded $T \to M$ link, and Stage~2 traverses the
$T$-adjusted $M \to Y$ link.  Together they reconstruct the total causal
effect $T \to Y$ without ever directly comparing the treatment groups on
$Y$, which is where the confounding by $U$ would otherwise enter.

\begin{example}{Smoking, Tar, and Cancer}
\label{w3:ex:smoking}
The classic front-door example \citep{pearl1995causal} is:
$T$ = smoking, $M$ = tar deposits in lungs, $Y$ = lung cancer,
$U$ = genetic predisposition (unobserved).
The graph is $T \to M \to Y$ with $U \to T$ and $U \to Y$.
There is no direct $T\to Y$ path (all effect of smoking operates through
tar).

Conditions~1--3 are satisfied: all causal paths from $T$ to $Y$ pass
through $M$; no back-door path exists from $T$ to $M$ (there is no arrow
from $U$ into $M$); and the only back-door path from $M$ to $Y$ is
$M \leftarrow T \leftarrow U \to Y$, which is blocked by conditioning on $T$.

The front-door formula gives:
\[
  P(\text{cancer} \mid \doop(T{=}\text{smoke}))
  = \int\!\left[\int P(\text{cancer} \mid t', m)\,dP(t')\right]
    dP(m \mid \text{smoke}),
\]
which is entirely computable from the observed joint distribution of
$(T, M, Y)$, without any data on $U$.  For the discrete case this becomes
a double sum over $(m, t')$.
\end{example}

\begin{example}{Front-Door Formula: Binary Numerical Computation}
\label{w3:ex:numerical}
Let $T, M, Y \in \{0, 1\}$ with $U$ unobserved and the graph
$T \to M \to Y$, $U \to T$, $U \to Y$.  The observed joint distribution
$P(T, M, Y)$ is given in the following population table.

\medskip
\begin{center}
\renewcommand{\arraystretch}{1.3}
\begin{tabular}{@{\hspace{8pt}}ccc@{\hspace{16pt}}c@{\hspace{8pt}}}
\toprule
$T$ & $M$ & $Y$ & $P(T,M,Y)$ \\
\midrule
0 & 0 & 0 & 0.360 \\
0 & 0 & 1 & 0.040 \\
0 & 1 & 0 & 0.050 \\
0 & 1 & 1 & 0.050 \\
1 & 0 & 0 & 0.070 \\
1 & 0 & 1 & 0.030 \\
1 & 1 & 0 & 0.120 \\
1 & 1 & 1 & 0.280 \\
\bottomrule
\end{tabular}
\end{center}

\medskip\noindent
All quantities needed for the front-door formula can be read off from this
table by marginalization.  For example,
$P(T{=}1) = 0.070 + 0.030 + 0.120 + 0.280 = 0.5$,
and
$P(M{=}1 \mid T{=}1) = (0.120 + 0.280)/(0.5) = 0.8$.

\medskip\noindent
We compute $P(Y{=}1 \mid \doop(T{=}t))$ for $t \in \{0,1\}$ using the
front-door formula.  The inner sum (Stage~2) does not depend on the
intervention value $t$, so compute it once for each value of $M$:
\[
  S_0 \;=\; \sum_{t'} P(Y{=}1 \mid t', M{=}0)\,P(t')
         = 0.1 \times 0.5 + 0.3 \times 0.5 = 0.20,
\]
\[
  S_1 \;=\; \sum_{t'} P(Y{=}1 \mid t', M{=}1)\,P(t')
         = 0.5 \times 0.5 + 0.7 \times 0.5 = 0.60.
\]
Now apply the outer sum (Stage~1) for each intervention value:
\begin{align*}
  P(Y{=}1 \mid \doop(T{=}1))
    &= P(M{=}0 \mid T{=}1)\,S_0 + P(M{=}1 \mid T{=}1)\,S_1 \\
    &= 0.2 \times 0.20 \;+\; 0.8 \times 0.60 \;=\; 0.52, \\[4pt]
  P(Y{=}1 \mid \doop(T{=}0))
    &= P(M{=}0 \mid T{=}0)\,S_0 + P(M{=}1 \mid T{=}0)\,S_1 \\
    &= 0.8 \times 0.20 \;+\; 0.2 \times 0.60 \;=\; 0.28.
\end{align*}
The average causal effect is $0.52 - 0.28 = 0.24$.

\paragraph{Comparison with the naive approach.}
A naive analyst ignoring the unobserved confounder $U$ would estimate the
effect of $T$ on $Y$ by the observational contrast
$P(Y{=}1 \mid T{=}1) - P(Y{=}1 \mid T{=}0)$, read directly from the
population table:
\begin{align*}
  P(Y{=}1 \mid T{=}1)
    &= \frac{0.030 + 0.280}{0.5} = 0.62, \\[4pt]
  P(Y{=}1 \mid T{=}0)
    &= \frac{0.040 + 0.050}{0.5} = 0.18.
\end{align*}
The naive estimate of the causal effect is $0.62 - 0.18 = 0.44$, which is
\emph{almost twice} the true causal effect of $0.24$.  The bias arises
because $U$ simultaneously increases $T$ and $Y$: high-$U$ individuals
tend to have $T{=}1$ and $Y{=}1$ for reasons unrelated to the $T\to M\to Y$
pathway, inflating the observational association.  The front-door formula
removes this bias by routing the effect estimate through the unconfounded
$T\to M$ link and the $T$-adjusted $M\to Y$ link, never directly comparing
the $T{=}1$ and $T{=}0$ groups on $Y$.

Note how Stage~2 (the inner sum $S_m$) averages $P(Y\mid t', m)$ over the
\emph{marginal} distribution $P(t')$, not the treatment-conditional
distribution $P(t' \mid m)$.  This re-weighting is precisely what removes
the confounding introduced by $U$: the Stage~2 quantity $S_m$ is the
back-door-adjusted effect of $M$ on $Y$, with $T$ serving as the
adjustment variable.
\end{example}

\begin{warn}{Comparing Back-Door and Front-Door}
The two criteria are \emph{not} nested.  Neither implies the other.
\begin{itemize}
  \item The back-door criterion requires all confounders to be observed
        (or blocked by observed variables), but places no restriction on
        mediators.
  \item The front-door criterion requires an observed mediator intercepting
        all causal paths, but the confounder $U$ can be entirely
        unobserved.
\end{itemize}
There exist graphs where neither criterion alone applies, but where the
full do-calculus (or the ID algorithm) can still identify the causal
effect.
\end{warn}

% ==============================================================
\section{The Three Rules of Do-Calculus}
\label{w3:sec:rules}
% ==============================================================

Having seen two concrete identification strategies, we now develop the
general algebraic machinery that underlies both.  \citet{pearl1995causal}
introduced the three rules of do-calculus as a system for manipulating
expressions involving the do-operator: each rule licenses a transformation
that removes or exchanges terms, bringing the expression one step closer
to a purely observational formula.  The back-door and front-door formulas
of Sections~\ref{w3:sec:backdoor} and~\ref{w3:sec:frontdoor} are each
derivable as three-step sequences of these rules, as shown in
Section~\ref{w3:sec:proofs}.  Whether the three rules are
\emph{sufficient} to identify every identifiable causal effect is a deeper
question; Section~\ref{w3:sec:completeness} presents the completeness
theorem of \citet{shpitser2006identification}, which answers it
affirmatively.

\subsection{Statement of the Rules}
\label{w3:sec:rules-statement}
Recall from Section~\ref{w2:sec:markov} of Chapter~2 that a distribution
$P$ is \emph{Markov with respect to $\Gcal$} if every d-separation statement in $\Gcal$
implies the corresponding conditional independence in $P$:
\[
  (X \indep Y \mid \mathbf{S})_{\Gcal}
  \;\Longrightarrow\;
  X \indep Y \mid \mathbf{S} \text{ under } P.
\]
Equivalently, $P$ admits the Markov factorization
$P(v_1,\ldots,v_n) = \prod_i P(v_i \mid \Pa(v_i))$.

\paragraph{How to read the graphical conditions in Theorem~\ref{w3:thm:docalculus}.}
Each rule has the form \emph{``expression~A = expression~B if
$(Y \indep Z \mid X, W)_{\Gcal'}$''}, where $\Gcal'$ is a specific
intervention graph built from $\Gcal$ by two deletions.
The subscript encodes both deletions in compressed notation.
The table below gives the recipe for each subscript that appears.

\medskip
\begin{center}
\renewcommand{\arraystretch}{1.45}
\begin{tabular}{@{\hspace{6pt}}l@{\hspace{14pt}}l@{\hspace{14pt}}l@{\hspace{6pt}}}
\toprule
\textbf{Subscript} & \textbf{Step 1: delete} & \textbf{Step 2: delete} \\
\midrule
$\Gcal_{\overline{X}\,\underline{Z}}$
  & all arrows \emph{into} $X$
  & all arrows \emph{out of} $Z$ \\
$\Gcal_{\overline{X}\,\overline{Z(W)}}$
  & all arrows \emph{into} $X$
  & arrows \emph{into} $Z$ except those from $W$ \\
$\Gcal_{\overline{X}\,\overline{Z(\overline{W})}}$
  & all arrows \emph{into} $X$
  & arrows \emph{into} $Z$ that are not from ancestors of $W$
    in $\Gcal_{\overline{X}}$ \\
\bottomrule
\end{tabular}
\end{center}

\medskip\noindent
Having formed $\Gcal'$ by these two deletions, the condition
$(Y \indep Z \mid X, W)_{\Gcal'}$ is then checked by the same
d-separation algorithm used in Chapter~2: find all paths between $Y$ and
$Z$, and verify that every path is blocked by $\{X, W\}$ in $\Gcal'$.

\begin{thm}{The Three Rules of Do-Calculus \citep{pearl1995causal}}
\label{w3:thm:docalculus}
Let $\Gcal$ be a DAG over variables $\mathcal{V}$, and let $X, Y, Z, W
\subseteq \mathcal{V}$ be disjoint sets.  The rules hold for any
distribution $P$ Markov with respect to $\Gcal$.

\medskip
\noindent\textbf{Rule 1 (Insertion/Deletion of Observations).}
\[
  P(y \mid \doop(x), z, w) \;=\; P(y \mid \doop(x), w)
  \quad\text{if}\quad
  (Y \indep Z \mid X, W)_{\Gcal_{\overline{X}\,\underline{Z}}}.
\]
If $Y$ and $Z$ are d-separated in the graph with arrows into $X$ deleted
\emph{and} arrows out of $Z$ deleted, then conditioning on $Z$ has no
effect on the distribution of $Y$ (given $\doop(x)$ and $W$).

\medskip
\noindent\textbf{Rule 2 (Action/Observation Exchange).}
\[
  P(y \mid \doop(x), \doop(z), w) \;=\; P(y \mid \doop(x), z, w)
  \quad\text{if}\quad
  (Y \indep Z \mid X, W)_{\Gcal_{\overline{X}\,\overline{Z(W)}}}.
\]
When the condition holds, intervening on $Z$ is equivalent to conditioning
on $Z$ (within the scope of $\doop(x)$ and $W$).

\medskip
\noindent\textbf{Rule 3 (Insertion/Deletion of Actions).}
\[
  P(y \mid \doop(x), \doop(z), w) \;=\; P(y \mid \doop(x), w)
  \quad\text{if}\quad
  (Y \indep Z \mid X, W)_{\Gcal_{\overline{X}\,\overline{Z(\overline{W})}}}.
\]
When the condition holds, the intervention $\doop(z)$ can be deleted
entirely.
\end{thm}

\begin{remark}
Each rule is a \emph{conditional independence statement in a specific
intervention graph}.  The three rules cover all ways of adding or removing
terms from a do-expression: Rule~1 handles observations; Rule~2 exchanges
actions for observations; Rule~3 deletes actions.  Any sequence of such
operations that converts a do-expression into a purely observational
expression is an \emph{identification proof}.
\end{remark}

\begin{center}
\renewcommand{\arraystretch}{1.35}
\begin{tabular}{@{\hspace{6pt}}c@{\hspace{12pt}}|@{\hspace{12pt}}l@{\hspace{6pt}}}
\textbf{Rule} & \textbf{Interpretation} \\
\hline
Rule 1 & Insert or delete \emph{observations} from a conditioning set \\
Rule 2 & Exchange an \emph{intervention} $\doop(z)$ for an \emph{observation} $z$ \\
Rule 3 & Insert or delete \emph{interventions} from an expression \\
\end{tabular}
\end{center}

\subsection{Intuition for Each Rule}
\label{w3:sec:rules-intuition}
\paragraph{Rule 1 --- Adding/Removing Observations.}
The condition $(Y \indep Z \mid X, W)_{\Gcal_{\overline{X}\,\underline{Z}}}$
says: in the post-$\doop(x)$ world, if we further deactivate $Z$'s causal
influence (delete its outgoing arrows), $Y$ and $Z$ become independent given
$W$.  This means $Z$ carries no information about $Y$ that is not already
in $W$, so it can be freely inserted or removed from the conditioning set.

\begin{example}{Rule 1 --- Deleting a Redundant Observation}
\label{w3:ex:rule1}
Graph: $Z \to T \to Y$, no other edges.  Suppose we intervene on $T$.

\begin{center}
\begin{tikzpicture}
  \node[node] (Z) at (0,0)   {$Z$};
  \node[node] (T) at (1.8,0) {$T$};
  \node[node] (Y) at (3.6,0) {$Y$};
  \draw[edge] (Z) -- (T);
  \draw[edge] (T) -- (Y);
\end{tikzpicture}
\end{center}

\noindent
In $\Gcal_{\overline{T}\,\underline{Z}}$: delete arrows into $T$ (removes
$Z \to T$), then delete arrows out of $Z$ (nothing remains).  Node $Z$ is
now isolated, so $(Y \indep Z \mid T)$ holds trivially.  Rule~1 gives:
\[
  P(y \mid \doop(t), z) \;=\; P(y \mid \doop(t)).
\]
\textit{Interpretation}: the intervention $\doop(t)$ severs $Z$'s only
route to $Y$.  Once we fix $T$ externally, knowing $Z$ (the former cause of
$T$) carries no additional information about $Y$.
\end{example}

\paragraph{Rule 2 --- Swapping Action for Observation.}
This is the most frequently used rule.  The condition
$(Y \indep Z \mid X, W)_{\Gcal_{\overline{X}\,\overline{Z(W)}}}$ says: in
a graph where we delete all back-door paths into $Z$ (except those from
$W$), $Y$ and $Z$ are d-separated given $X$ and $W$.  Intuitively, if
the only paths from $Z$ to $Y$ are causal (forward) paths, then intervening
on $Z$ and conditioning on $Z$ have the same effect on $Y$ --- there is no
confounding between $Z$ and $Y$.

\begin{example}{Rule 2 --- Exchanging Intervention for Observation}
\label{w3:ex:rule2}
Graph: $Z \to T \to Y$, no other edges.  Suppose we have already intervened
on $T$ and wish to exchange the intervention $\doop(z)$ for the observation
$z$ in the expression $P(y \mid \doop(z), t)$.

\begin{center}
\begin{tikzpicture}
  \node[node] (Z) at (0,0)   {$Z$};
  \node[node] (T) at (1.8,0) {$T$};
  \node[node] (Y) at (3.6,0) {$Y$};
  \draw[edge] (Z) -- (T);
  \draw[edge] (T) -- (Y);
\end{tikzpicture}
\end{center}

\noindent
Here $X = \varnothing$, $Z = Z$, $W = \{T\}$.  Form
$\Gcal_{\overline{Z(T)}}$: delete arrows into $Z$ except those from $T$.
Since $Z$ has no incoming arrows, $\Gcal_{\overline{Z(T)}} = \Gcal$.
Check $(Y \indep Z \mid T)_{\Gcal}$: the only path from $Z$ to $Y$ is
$Z \to T \to Y$, which is a chain blocked by conditioning on $T$. \checkmark

Rule~2 gives:
\[
  P(y \mid \doop(z), t) \;=\; P(y \mid z, t).
\]
\textit{Interpretation}: $T$ mediates the entire influence of $Z$ on $Y$.
Once $T$ is observed, the distribution of $Y$ is the same whether $Z$ was
set by an external intervention or merely observed to equal $z$ --- $T$
screens off $Z$ from $Y$ completely.
\end{example}

\paragraph{Rule 3 --- Deleting an Intervention.}
The condition says that $Z$ has no causal path to $Y$ that is not blocked by
$X$ or $W$ in the appropriate intervention graph.  Under this condition,
setting $Z$ to any value has no effect on $Y$, so $\doop(z)$ can be dropped.

\begin{example}{Rule 3 --- Deleting an Irrelevant Intervention}
\label{w3:ex:rule3}
Same graph as Example~\ref{w3:ex:rule1}: $Z \to T \to Y$, no other edges.
Now suppose we intervene on \emph{both} $Z$ and $T$.

\begin{center}
\begin{tikzpicture}
  \node[node] (Z) at (0,0)   {$Z$};
  \node[node] (T) at (1.8,0) {$T$};
  \node[node] (Y) at (3.6,0) {$Y$};
  \draw[edge] (Z) -- (T);
  \draw[edge] (T) -- (Y);
\end{tikzpicture}
\end{center}

\noindent
In $\Gcal_{\overline{T}}$ (delete arrows into $T$: removes $Z \to T$),
node $Z$ has no causal path to $Y$.  Rule~3 gives:
\[
  P(y \mid \doop(t), \doop(z)) \;=\; P(y \mid \doop(t)).
\]
\textit{Interpretation}: with $T$ already fixed by intervention, $Z$'s
only route to $Y$ (through $T$) is blocked.  Setting $Z$ to any value
cannot change $Y$, so the intervention $\doop(z)$ is superfluous and can
be deleted.

\textit{Contrast with Rule~1}: Example~\ref{w3:ex:rule1} showed that the
\emph{observation} $z$ can also be deleted from $P(y \mid \doop(t), z)$
for the same reason.  Rule~1 deletes a redundant \emph{observation};
Rule~3 deletes a redundant \emph{intervention}.  Both apply here because
$\doop(t)$ renders $Z$ causally inert.
\end{example}

\subsection{Proofs of the Identification Formulas}
\label{w3:sec:proofs}
With the three rules now in hand, we can prove the two identification
formulas stated in Sections~\ref{w3:sec:backdoor}
and~\ref{w3:sec:frontdoor}.

\begin{tcolorbox}[defstyle, title={Proof of Theorem~\ref{w3:thm:backdoor} (Back-Door Formula) via Rules~2 and~3}]

\noindent\textit{Graph structure}: $T \to Y$, unobserved $U \to T$,
$U \to Y$, and observed $\mathbf{S}$ blocking the back-door path
$T \leftarrow U \to Y$.

\begin{center}
\begin{tikzpicture}[node distance=1.6cm]
  \node[node]  (T) at (0, 0)     {$T$};
  \node[node]  (Y) at (3.0, 0)   {$Y$};
  \node[unode] (U) at (1.5, 1.2) {$U$};
  \node[node]  (Z) at (1.5,-1.2) {$\mathbf{S}$};
  \draw[edge]  (T) -- (Y);
  \draw[dedge] (U) -- (T);
  \draw[dedge] (U) -- (Y);
  \draw[edge]  (Z) -- (T);
  \draw[edge]  (Z) -- (Y);
  \node[font=\scriptsize\color{warnred}] at (1.5, 0.35)
    {\textit{back-door path}};
\end{tikzpicture}
\end{center}

\noindent The back-door path $T \leftarrow U \to Y$ is shown dashed.
$\mathbf{S}$ blocks this path.  The proof manipulates $P(y \mid \doop(t))$
in three steps.

\medskip
\noindent\textbf{Step 1: Introduce $\mathbf{S}$ by the law of total
probability.}
\[
  P(y \mid \doop(t))
  = \int P(y \mid \doop(t),\;\mathbf{s})\,dP(\mathbf{s} \mid \doop(t)).
\]

\noindent\textbf{Step 2: Simplify $P(\mathbf{s} \mid \doop(t))$.}
Because $\mathbf{S}$ contains no descendants of $T$ (condition~1),
intervening on $T$ cannot causally affect any node in $\mathbf{S}$.  In
the intervention graph $\Gcal_{\overline{T}}$, all arrows into $T$ have been
deleted, so $T$ has no parents and no causal path from $T$ reaches
$\mathbf{S}$.  By Rule~3:
\[
  P(\mathbf{s} \mid \doop(t)) = P(\mathbf{s}).
\]

\noindent\textbf{Step 3: Convert $P(y \mid \doop(t), \mathbf{s})$ to an
observation.}
Condition~2 guarantees that $\mathbf{S}$ blocks all back-door paths from
$T$ to $Y$.  In $\Gcal_{\overline{T}}$, the only remaining paths from $T$
to $Y$ are causal paths.  Rule~2 then licenses the exchange:
\[
  P(y \mid \doop(t), \mathbf{s}) = P(y \mid t, \mathbf{s}).
\]

\noindent\textbf{Combining:}
\[
  P(y \mid \doop(t))
  = \int P(y \mid t, \mathbf{s})\,dP(\mathbf{s}).
  \hfill\square
\]
\end{tcolorbox}

\begin{tcolorbox}[defstyle, title={Proof of Theorem~\ref{w3:thm:frontdoor} (Front-Door Formula) via Three Applications of the Rules}]

\noindent\textit{Graph structure}: $T \to M \to Y$, unobserved $U \to T$,
$U \to Y$, no direct $T \to Y$ edge.

\begin{center}
\begin{tikzpicture}[node distance=1.6cm]
  \node[node]  (T) at (0, 0)     {$T$};
  \node[node]  (M) at (2.2, 0)   {$M$};
  \node[node]  (Y) at (4.4, 0)   {$Y$};
  \node[unode] (U) at (2.2, 1.4) {$U$};
  \draw[edge]  (T) -- (M) node[midway, below,
    font=\scriptsize\itshape\color{darkgrey}] {Stage 1};
  \draw[edge]  (M) -- (Y) node[midway, below,
    font=\scriptsize\itshape\color{darkgrey}] {Stage 2};
  \draw[dedge] (U) -- (T);
  \draw[dedge] (U) -- (Y);
\end{tikzpicture}
\end{center}

\noindent Stage~1 ($T \to M$) is unconfounded by condition~2.
Stage~2 ($M \to Y$) is confounded by $U$ via $M \leftarrow T \leftarrow U
\to Y$, but conditioning on $T$ blocks this path (condition~3).

\medskip
\noindent\textbf{Step 1: Introduce $M$ by the law of total probability.}
\[
  P(y \mid \doop(t)) = \int P(y \mid \doop(t),\;m)\,dP(m \mid \doop(t)).
\]

\noindent\textbf{Step 2: Simplify $P(m \mid \doop(t))$.}
In $\Gcal_{\overline{T}}$, there are no back-door paths from $T$ to $M$
by front-door condition~2.  Rule~2 applies:
\[
  P(m \mid \doop(t)) = P(m \mid t).
\]

\noindent\textbf{Step 3: Simplify $P(y \mid \doop(t), m)$.}

\medskip
\noindent\textit{Step 3a}: $P(y \mid \doop(t), m) = P(y \mid \doop(m))$.

In $\Gcal_{\overline{T}}$, the arrow $U \to T$ is deleted, so $U$ and $T$
become independent.  Since there is no arrow from $U$ into $M$ (condition~2),
$U$ is also independent of $M$ in $\Gcal_{\overline{T}}$.  Therefore:
\[
  P(y \mid \doop(t), m) \;=\; \int P(y \mid m, u)\,dP(u) \;=\; P(y \mid \doop(m)).
\]
This also follows from Rule~3: in $\Gcal_{\overline{M}\,\overline{T}}$,
$T$ is isolated and $(Y \indep T)$ holds, licensing
$P(y \mid \doop(t), \doop(m)) = P(y \mid \doop(m))$.

\medskip
\noindent\textit{Step 3b}: $P(y \mid \doop(m)) = \int P(y \mid m, t')\,dP(t')$.

By front-door condition~3, $\{T\}$ satisfies the back-door criterion for
the effect of $M$ on $Y$, and Theorem~\ref{w3:thm:backdoor} gives:
\[
  P(y \mid \doop(m)) \;=\; \int P(y \mid m, t')\,dP(t').
\]

\noindent Combining Steps 3a and 3b:
\[
  P(y \mid \doop(t), m) \;=\; \int P(y \mid m, t')\,dP(t').
\]

\noindent\textbf{Assembling:}
\[
  P(y \mid \doop(t))
  = \int\!\left[\int P(y \mid t',\,m)\,dP(t')\right]dP(m \mid t).
  \hfill\square
\]
\end{tcolorbox}

% ==============================================================
\section{The Do-Calculus is Complete}
\label{w3:sec:completeness}
% ==============================================================

A natural question after seeing the three rules is: are they enough?
Could there be a graph where the causal effect is identifiable in principle,
but none of the rules --- applied in any sequence --- can derive a purely
observational expression?  \citet{shpitser2006identification} answered this
definitively.

Throughout Sections~\ref{w3:sec:mutilated}--\ref{w3:sec:rules} we have
represented latent common causes as
explicit unobserved nodes (e.g.\ $U \to T$ and $U \to Y$).  The
completeness theorem uses a more general graph class that encodes latent
common causes compactly as \emph{bidirected edges}: a bidirected edge
$X \leftrightarrow Y$ stands for an unobserved variable $U$ with $U \to X$
and $U \to Y$.  A DAG augmented with such bidirected edges is called a
\emph{semi-Markovian DAG}.  The IV DAG and front-door graphs of this course
are both special cases.

\begin{thm}{Completeness of the Do-Calculus \citep{shpitser2006identification}}
\label{w3:thm:completeness}
Let $\Gcal$ be a semi-Markovian DAG.  A causal quantity
$P(y \mid \doop(t))$ is identifiable from $P$ relative to $\Gcal$ if and
only if the do-calculus can derive a purely observational expression for it.
\end{thm}

The proof is beyond the scope of this chapter; the sketch below is
provided for reference and independent reading.

\begin{tcolorbox}[defstyle, title={Proof Sketch and the ID Algorithm}]
\citet{shpitser2006identification} introduced the \emph{ID algorithm}
(Identification algorithm): a procedure that takes as input a
semi-Markovian DAG $\Gcal$, a target intervention distribution
$P(y \mid \doop(t))$, and the observed joint distribution $P$, and either
returns an explicit observational formula for $P(y \mid \doop(t))$ or
reports that the quantity is not identifiable from $P$ given $\Gcal$.

\textbf{Sufficiency} (do-calculus identifies whenever identification is
possible): proved constructively via the ID algorithm.  The key structural
concept is the \emph{c-component} (confounded component): a maximal set of
nodes connected by bidirected edges, representing a group of variables that
share latent common causes.  For a single node $V$ with no bidirected
edges, its c-component is $\{V\}$ itself.

The ID algorithm proceeds recursively as follows.
\begin{enumerate}
  \item \textit{Simplification}: remove from $\Gcal$ any node that is not
        an ancestor of $Y$ in $\Gcal_{\overline{T}}$.  Such nodes cannot
        affect $P(y \mid \doop(t))$ and can be marginalized out.
  \item \textit{Decomposition by c-components}: factorise
        $P(\mathbf{v} \setminus \mathbf{t} \mid \doop(t))$ over the
        c-components of $\Gcal_{\overline{T}}$.  Each factor depends only
        on the variables in one c-component and their ancestors, reducing
        the problem to smaller subproblems.
  \item \textit{Recursive identification}: attempt to identify each
        c-component factor from $P$ using the do-calculus rules and the
        Markov factorization.  If every factor is successfully reduced to
        an observational expression, assemble the results to obtain the
        final formula for $P(y \mid \doop(t))$.
  \item \textit{Failure}: if any subproblem cannot be reduced --- because
        the c-component of $Y$ is not identifiable given the surrounding
        structure --- the algorithm returns \textsc{fail}.
\end{enumerate}
When the algorithm succeeds, it has produced an explicit identification
formula, demonstrating that the do-calculus is sufficient.

\textbf{Necessity} (do-calculus cannot identify non-identifiable
quantities): proved by showing that whenever the ID algorithm fails, the
causal effect is genuinely non-identifiable --- not merely hard to compute.
The key graphical obstruction is the \emph{hedge}.

\begin{quote}
\textit{A \emph{hedge} for $P(y \mid \doop(t))$ in $\Gcal$ is a pair of
subgraphs $(F, F')$ where $F$ and $F'$ are both rooted at the same set of
nodes, $F$ is a subgraph of $\Gcal$ containing a back-door path from $T$
to $Y$, and $F'$ is a subgraph of $\Gcal_{\overline{T}}$ in which $Y$
remains connected to a latent common cause.  Intuitively, the hedge
encodes a ``loop'' of confounding that the do-calculus cannot break.}
\end{quote}

\citet{shpitser2006identification} prove that the ID algorithm fails if and
only if $\Gcal$ contains a hedge for $P(y \mid \doop(t))$.  When a hedge
exists, one can explicitly construct two semi-Markovian models
$\mathcal{M}_1$ and $\mathcal{M}_2$ that:
\begin{itemize}
  \item induce the same observed joint distribution $P$, so no data can
        distinguish them; but
  \item assign different values to $P(y \mid \doop(t))$, so the
        interventional distribution is not a function of $P$ alone.
\end{itemize}
The construction builds $\mathcal{M}_1$ and $\mathcal{M}_2$ by choosing
different structural equations for the latent variables along the hedge,
keeping all observed conditional distributions identical.  Since the two
models are observationally indistinguishable yet causally distinct, no
method --- including the do-calculus --- can identify $P(y \mid \doop(t))$
from $P$ alone.  This establishes necessity.

Combining sufficiency and necessity: the do-calculus succeeds if and only
if no hedge exists, which is if and only if $P(y \mid \doop(t))$ is
identifiable. $\square$
\end{tcolorbox}

\begin{remark}
Theorem~\ref{w3:thm:completeness} has a crucial practical implication:
\emph{if the ID algorithm reports non-identification, then no estimation
method --- however clever --- can recover the causal effect from
observational data alone}, given the assumed graph structure.
Non-identification is not a limitation of a particular technique; it is a
fundamental property of the causal model.  Three remedies exist: impose
additional structural assumptions (e.g., linearity or monotonicity) that
go beyond the graph; collect additional data that break the obstruction
(e.g., an instrument or a randomized experiment); or target a different,
identified causal quantity instead (e.g., partial identification bounds,
or the effect among the treated rather than the average treatment effect).
\end{remark}

% ==============================================================
\section{The Big Picture}
\label{w3:sec:bigpicture}
% ==============================================================

This chapter completes the \emph{theoretical} identification machinery of
the course.  The chain of reasoning begun in Chapter~1 (the
interventional/observational distinction) and formalised in Chapter~2
(d-separation) now terminates in a concrete identification formula.
Chapters~4--7 then apply this machinery to specific research designs:

\bigskip
\begin{center}
\begin{tikzpicture}[
    box/.style = {
      rectangle, rounded corners=6pt,
      draw=isublue, fill=defbg, thick,
      text width=7.2cm, align=center,
      minimum height=1.1cm,
      font=\bfseries\color{isublue}
    },
    tool/.style = {
      rectangle, rounded corners=3pt,
      draw=accent, fill=white, thin,
      text width=4.8cm, align=center,
      font=\small\itshape\color{accent}
    },
    arr/.style = {-{Stealth[length=7pt,width=5pt]},
                  line width=1.6pt, color=accent},
    node distance = 0.55cm
  ]
  \node[box] (G)  {Graph Structure ($\Gcal$)};
  \node[box, below=1.2cm of G]  (D)
    {Mutilated Graphs ($\Gcal_{\overline{T}},\, \Gcal_{\underline{T}}$)};
  \node[box, below=1.2cm of D]  (R)
    {Do-Calculus Rules (1, 2, 3)};
  \node[box, below=1.2cm of R]  (ID)
    {Identification Formula $P(y\mid\doop(t)) = \Psi(P_{\mathrm{obs}})$};
  \node[box, below=1.2cm of ID] (W46)
    {Research Designs: Potential Outcomes, Mediation, Propensity Scores, IV
     (Chapters~4--8)};
  \node[box, below=1.2cm of W46] (ES)
    {Estimation: $\hat\Psi_n$ (Chapters~9--12)};

  \draw[arr] (G)   -- (D);
  \draw[arr] (D)   -- (R);
  \draw[arr] (R)   -- (ID);
  \draw[arr] (ID)  -- (W46);
  \draw[arr] (W46) -- (ES);

  \node[tool, right=1.1cm of G]
    {Causal assumptions encoded as directed edges};
  \node[tool, right=1.1cm of D]
    {$\Gcal_{\overline{T}}$ models intervention;
     $\Gcal_{\underline{T}}$ is a technical device for Rule~1};
  \node[tool, right=1.1cm of R]
    {Three rules license algebraic transformations};
  \node[tool, right=1.1cm of ID]
    {Back-door, front-door, or ID algorithm};
  \node[tool, right=1.1cm of W46]
    {Apply identification framework to concrete designs};
  \node[tool, right=1.1cm of ES]
    {GMM / EIF / doubly robust (Part~III)};

  \foreach \n in {G, D, R, ID, W46, ES}
    \draw[dotted, color=accent!50, thin] (\n.east) -- ++(1.1cm,0);
\end{tikzpicture}
\end{center}

\medskip
Chapter~4 introduces the potential outcomes framework and its connection
to the do-calculus, providing the second language for causal inference.
Chapters~5--6 apply these tools to randomization and back-door adjustment,
and to propensity score methods.  Chapter~7 develops instrumental variables
as the primary strategy when unmeasured confounding rules out the back-door
route.  Chapter~8 then studies mediation analysis, decomposing total effects
into direct and indirect components.
The completeness theorem (Section~\ref{w3:sec:completeness}) closes the
loop: every identifiable
causal effect can be found by the do-calculus, and every non-identifiable
one cannot be found by any method.  Part~III (Chapter~9 and beyond) then
turns to estimation.

\medskip

% ==============================================================
\section{Summary}
\label{w3:sec:summary}
% ==============================================================

\begin{enumerate}
  \item Two graph operations underpin the do-calculus:
        $\Gcal_{\overline{X}}$ (delete arrows \emph{into} $X$, modeling
        the intervention $\doop(x)$) and $\Gcal_{\underline{X}}$ (delete
        arrows \emph{out of} $X$, used as a technical device to check the
        d-separation condition for Rule~1 --- it does not model
        conditioning on $X$).

  \item The back-door criterion identifies $P(y \mid \doop(t))$ when an
        observed set $\mathbf{S}$ blocks all back-door paths and contains
        no descendant of $T$.  The general formula is
        $\int P(y \mid t, \mathbf{s})\,dP(\mathbf{s})$,
        equivalently $\E_{\mathbf{S}}\!\left[\E[h(Y)\mid T{=}t,\mathbf{S}]\right]$
        in iterated-expectation form (standardization; the g-formula for
        $K=1$).

  \item The front-door criterion identifies $P(y \mid \doop(t))$ via an
        observed mediator $M$ intercepting all causal paths from $T$ to
        $Y$, even when the confounder $U$ is entirely unobserved, provided
        the three front-door conditions hold.

  \item The do-calculus has three rules, each licensed by a d-separation
        condition in an appropriate intervention graph.  Rule~1 adds/removes
        observations; Rule~2 exchanges actions for observations; Rule~3
        deletes actions.  The back-door and front-door formulas are each
        derivable as three-step sequences of these rules
        (Section~\ref{w3:sec:proofs}).

  \item The do-calculus is complete \citep{shpitser2006identification}: a
        causal effect is identified from observational data if and only if
        the do-calculus can derive a purely observational expression for it.
        Non-identifiability is a fundamental obstacle, not a technical one.
\end{enumerate}

% ==============================================================
\newpage
\section*{Problems}
\addcontentsline{toc}{section}{Problems}
% ==============================================================

\begin{enumerate}[label=\textbf{\arabic*.}, leftmargin=*, itemsep=10pt]

  \item \textbf{Warm-up: graph operations.}
  Let $\Gcal$ be the DAG with edges $Z \to T$, $T \to Y$, $U \to T$,
  $U \to Y$ (the IV DAG, with $U$ unobserved).
  \begin{enumerate}[label=(\alph*)]
    \item Draw $\Gcal_{\overline{T}}$ and identify all paths from $Z$ to
          $Y$ that remain open.
    \item Draw $\Gcal_{\underline{T}}$ and identify all paths from $Z$ to
          $Y$ that remain.  For each path, classify it as a causal chain,
          fork, or collider path, and state whether it is open or closed
          by default (without conditioning on any variable).
    \item In $\Gcal_{\overline{T}}$, is $Z \indep Y \mid \varnothing$?
          How does the absence of a direct edge $Z \to Y$ in the DAG
          relate to the exclusion restriction (the assumption that $Z$
          affects $Y$ only through $T$)?
  \end{enumerate}

  \item \textbf{Back-door practice.}
  Consider the DAG: $X \to T$, $X \to Y$, $T \to M$, $M \to Y$,
  $T \to Y$, where $X$ is observed.
  \begin{enumerate}[label=(\alph*)]
    \item List all back-door paths from $T$ to $Y$.
    \item Does $\{X\}$ satisfy the back-door criterion for the effect of
          $T$ on $Y$?  Write the resulting adjustment formula.
    \item Does $\{M\}$ satisfy the back-door criterion?  Explain why or
          why not.
    \item Does $\{X, M\}$ satisfy the back-door criterion?  Identify
          which condition of the criterion $M$ violates, and explain why
          this rules it out as an adjustment variable.
    \item Even if the formal criterion issue in (d) were somehow set aside,
          explain the substantive bias that arises from conditioning on a
          mediator $M$ that lies on a causal path $T \to M \to Y$.
  \end{enumerate}

  \item \textbf{Front-door identification.}
  Consider the DAG: $T \to M \to Y$, with $U \to T$ and $U \to Y$
  (unobserved $U$, no direct $T \to Y$ edge).
  \begin{enumerate}[label=(\alph*)]
    \item Verify that $M$ satisfies all three front-door conditions.
    \item Derive the front-door formula~\eqref{w3:eq:frontdoor} step by
          step, citing the do-calculus rule used at each step.
    \item Now add a direct edge $T \to Y$ to the graph.  Does $M$ still
          satisfy the front-door criterion?  Explain which condition fails.
  \end{enumerate}

  \item \textbf{Identification or non-identification?}
  For each of the following DAGs, determine whether
  $P(y \mid \doop(t))$ is identified \emph{non-parametrically} from
  observational data.  If identified, state the formula; if not, explain
  the structural obstruction.
  \begin{enumerate}[label=(\alph*)]
    \item $T \to Y$, $U \to T$, $U \to Y$, $U$ unobserved.
    \item $T \to Y$, $U \to T$, $U \to Y$, $U$ unobserved; with instrument
          $Z \to T$.
          (\textit{Hint}: in the non-parametric setting, an instrument
          alone does not identify the average causal effect
          $P(y\mid\doop(t))$ without additional structural assumptions
          such as linearity or monotonicity.  What the instrument
          \emph{does} identify --- the Wald estimand and the LATE --- is
          developed in Chapter~7.)
    \item $T \to M \to Y$, $U \to T$, $U \to Y$, $U \to M$, $U$
          unobserved.  (\textit{Hint}: first check whether the back-door
          criterion can be satisfied --- is there any observed variable
          that blocks the path $T \leftarrow U \to Y$?  Then check the
          front-door criterion: does $M$ satisfy condition~2?  Conclude
          whether the effect is identified by either strategy.)
    \item $T \to Y$, $X \to T$, $X \to Y$, $X$ observed.
  \end{enumerate}

\end{enumerate}

% ==============================================================
\newpage
\section*{Biographical Note: Judea Pearl}
\addcontentsline{toc}{section}{Biographical Note: Judea Pearl}
% ==============================================================

\begin{tcolorbox}[
  colback=defbg, colframe=isublue, arc=4pt,
  boxrule=0.8pt, left=8pt, right=8pt, top=6pt, bottom=6pt
]
\textbf{Judea Pearl} (born 1936, Tel Aviv) is a computer scientist and
philosopher whose work over four decades created the modern mathematical
foundations of causal inference and probabilistic reasoning.

Pearl trained as an electrical engineer, receiving his B.Sc.\ from the
Technion (Israel Institute of Technology) in 1960 and his Ph.D.\ from
the Polytechnic Institute of Brooklyn in 1965.  He joined the faculty of
UCLA in 1970, where he has remained ever since.  His early research was
in hardware design and adaptive systems, but in the late 1970s he turned
to artificial intelligence and the problem of reasoning under uncertainty.

His first landmark contribution was the \emph{Bayesian network} (1985): a
directed acyclic graph in which nodes represent random variables and edges
encode conditional independence structure via the Markov factorization.
Pearl's belief-propagation algorithm made exact inference in such networks
tractable, and Bayesian networks became the dominant framework for
probabilistic reasoning in AI throughout the 1980s and 1990s.

The second, and arguably deeper, contribution was the \emph{do-calculus}
(1993--1995).  Pearl recognized that Bayesian networks, despite their
power, could not distinguish association from causation: the same DAG that
encodes observational dependencies is also used to model interventions, but
the two require different algebra.  He introduced the $\doop(\cdot)$
operator to represent external interventions, the intervention-graph
semantics to compute their effects, and the three rules of do-calculus as a
complete axiomatic system for moving between interventional and
observational distributions.  His 2000 monograph
\textit{Causality: Models, Reasoning, and Inference} unified these ideas
into a comprehensive theory, covering identification, counterfactuals,
mediation, and the potential-outcomes framework as a special case.

In recognition of these contributions, Pearl was awarded the \textbf{ACM
Turing Award} in 2011 --- the highest honor in computer science,
often described as the ``Nobel Prize of computing'' --- ``for fundamental
contributions to artificial intelligence through the development of a
calculus for probabilistic and causal reasoning.''  He is also the
recipient of the Benjamin Franklin Medal in Computers and Cognitive Science
(2008) and numerous honorary doctorates.

Pearl has been a persistent advocate for what he calls the
\emph{ladder of causation}: a hierarchy distinguishing association
(seeing), intervention (doing), and counterfactual reasoning (imagining).
He argues that true intelligence --- human or artificial --- requires
climbing all three rungs, and that current machine learning systems, however
powerful, remain largely confined to the first.  His 2018 popular book
\textit{The Book of Why} (co-authored with Dana Mackenzie) brings these
ideas to a general audience.
\end{tcolorbox}

% ==============================================================
\newpage
\ifSubfilesClassLoaded{
  \bibliographystyle{plainnat}
  \bibliography{causal_inference}
}{}

\end{document}